
% Compile with XeLaTeX.
\documentclass[aspectratio=43,9pt]{beamer}

\usepackage{fontspec}
\usepackage{microtype}
\usepackage{mathtools}
\usepackage{booktabs}
\usepackage{tabularx}
\usepackage{array}
\usepackage{tikz}
\usepackage{colortbl}
\usepackage{ragged2e}
%\usepackage{enumitem}
\usepackage{graphicx}
\usepackage{adjustbox}
\usepackage{multirow}
\usepackage{xcolor}

\usetikzlibrary{arrows.meta,positioning,calc,fit,backgrounds}

\usepackage{hyperref}

\hypersetup{
    breaklinks=true
}

\Urlmuskip=0mu plus 1mu

% Fonts
\setmainfont{Liberation Sans}
\setsansfont{Liberation Sans}
\setmonofont{DejaVu Sans Mono}

% Theme
\usefonttheme{professionalfonts}
\usetheme[numbering=fraction,progressbar=frametitle]{metropolis}
\metroset{block=fill}
\setbeamertemplate{navigation symbols}{}
\setbeamertemplate{caption}{\raggedright\insertcaption\par}
\setbeamertemplate{itemize items}[bullet]
\setbeamertemplate{enumerate items}[default]
\setbeamertemplate{footline}{}

% Colors
\definecolor{Navy}{HTML}{0B1F3A}
\definecolor{Teal}{HTML}{0F766E}
\definecolor{Slate}{HTML}{475569}
\definecolor{SoftBg}{HTML}{F8FAFC}
\definecolor{Accent}{HTML}{2563EB}
\definecolor{WarmGray}{HTML}{64748B}
\definecolor{SymbiosisRed}{HTML}{C41230}

\setbeamercolor{normal text}{fg=Slate,bg=white}
\setbeamercolor{alerted text}{fg=Teal}
\setbeamercolor{frametitle}{fg=Navy,bg=white}
\setbeamercolor{progress bar}{fg=Teal,bg=SoftBg}
\setbeamercolor{block title}{fg=white,bg=Navy}
\setbeamercolor{block body}{fg=Slate,bg=SoftBg}
\setbeamercolor{block title alerted}{fg=white,bg=Teal}
\setbeamercolor{block body alerted}{fg=Slate,bg=SoftBg}
\setbeamercolor{block title example}{fg=white,bg=Accent}
\setbeamercolor{block body example}{fg=Slate,bg=SoftBg}
\setbeamercolor{title}{fg=Navy}
\setbeamercolor{subtitle}{fg=Slate}
\setbeamercolor{section in toc}{fg=Navy}
\setbeamercolor{structure}{fg=Teal}

% Spacing
%\setlist[itemize]{leftmargin=1.15em,itemsep=0.28em,topsep=0.15em}
%\setlength{\parskip}{0.25em}
%\setlength{\parindent}{0pt}
%\RaggedRight
%\sloppy
%\setlength{\leftmargini}{1.15em}

%\addtobeamertemplate{itemize/enumerate body begin}{%
%  \setlength{\itemsep}{0.3em}%
%  \setlength{\parskip}{0pt}%
%  \setlength{\parsep}{0pt}%
%  \RaggedRight%
%  \sloppy%
%}{}
%
%\setbeamertemplate{itemize items}[bullet]
%\setlength{\leftmargini}{1.15em}

\setbeamertemplate{itemize item}{\textbullet}
\setbeamertemplate{itemize subitem}{\textbullet}
\setbeamertemplate{itemize subsubitem}{\textbullet}

\newcommand{\kw}[1]{\textbf{#1}}

\newcommand{\figph}[3]{%
\begin{center}
\setlength{\fboxsep}{8pt}%
\fcolorbox{WarmGray}{SoftBg}{%
\begin{minipage}[c][#1][c]{\dimexpr #2-2\fboxsep-2\fboxrule\relax}
\centering
{\bfseries #3\par}
\vspace{0.35em}
{\footnotesize Insert figure, screenshot, chart, or diagram here.\par}
{\footnotesize Why it helps: anchors the slide in evidence and keeps the text sparse.\par}
\end{minipage}}%
\end{center}
}

\title{Progress Report}
\subtitle{26 Jun 2026 - 31 Jul 2026}
\author{Soham Kulkarni}
\date{31 July 2026}

\begin{document}

\begin{frame}[plain]
    \rmfamily % Switches to Times/Serif font for exact visual matching
    \centering
    \color{black}

    {\Large\bfseries B. Tech. Project AY 2026-27}\\[0.2cm]
    {\Large CA 1 Presentation}\\[0.1cm]

    {\large \textbf{Project ID:} INT2 \\ \textbf{Project}: Porting ROCm to Consumer AMD GPUs}
    
%    \vspace{0.2cm}

    % Logo Placeholder
    \begin{center}
        
\includegraphics[height=2.5cm]{SIT_Logo.png} % Replace 'example-image-a' with your logo file (e.g., logo.png)
    \end{center}


    {\large \textbf{Guides:} Dr. Shruti Sunnad \& Dr. Madhvi Saxena}\\[0.2cm]
    {\large \textbf{Industry Mentors}: \\ Rahul Poruri (CEO, FOSS United) \\\& \\Bhooshan Agalgatti (Sr. Director, Nvidia)}\\[0.2cm]
    {\large \textbf{Member}: Soham Kulkarni (23070126129)}

%\vspace{0.2cm}
\begin{color}{gray}
    \centering {\small Department of Artificial Intelligence \& Machine Learning}
\end{color}
\end{frame}

\begin{frame}{Project evolution and reporting window}
\begin{itemize}
        \item Current work targets \texttt{gfx90c}, with workflow validation planned on \texttt{gfx902}.
        \item Reporting period: \texttt{26-06-26} to \texttt{31-07-26} across two progress reports.
        \item Technical evidence is documented through eight project blog articles published on \href{https://netherquark.substack.com}{Substack}.
    \end{itemize}
    \begin{figure}[htbp]
        \includegraphics[width=\linewidth]{articles_matrix.png}
    \end{figure}
\end{frame}


\begin{frame}{Introduction}
\begin{itemize}
        \item Conducted as an employee of \textit{FOSS United Foundation}.
        \item ROCm is AMD's CUDA
        \item Primary target: \texttt{gfx90c} (Ryzen 5 5650U); validation target: \texttt{gfx902} (Ryzen 5 2200G).
        \item Goal: develop a reproducible workflow for adding ROCm support.
\end{itemize}

\vspace{0.5cm}
\begin{center}
\begin{figure}[htbp]
    
\includegraphics[width=0.75\linewidth]{AMD_Logo.png}
\end{figure}
\end{center}

\end{frame}


\begin{frame}{Problem statement}
\begin{itemize}
        \item Native HIP execution succeeds, but packaged ROCm libraries lack complete gfx90c support.
        \item Missing rocBLAS/Tensile artifacts prevent higher-level workloads such as PyTorch.
        \item Community spoofing enables evaluation but is not a production-grade solution.
        \item A reproducible recovery workflow for unsupported architectures is currently unavailable.
\end{itemize}
\end{frame}


\begin{frame}{Objectives of the project}
\begin{itemize}
        \item Recover functional ROCm support for gfx90c using reproducible methods.
        \item Benchmark every recovery stage against the native software stack \& against spoofing.
        \item Validate the methodology on gfx902.
        \item Generalise the workflow for other unsupported AMD architectures.
    \end{itemize}
\end{frame}

\begin{frame}{Project Timeline}
\vspace{-0.5cm}
\begin{itemize}
        \item Jun 2026: Verify feasibility of project
        \item Jul 2026: Port 1 HIP application from \texttt{gfx900} to \texttt{gfx90c}
        \item Aug 2026: Develop expertise in GCN Assembly, HIP, ROCm build system
        \item Sep 2026: Finish porting \texttt{rocm-examples}
        \item Oct 2026: Finish porting full ROCm stack
        \item Nov 2026: Finish porting PyTorch / Dockerised PyTorch
        \item Dec 2026: Establish a reproducible recovery workflow for \texttt{gfx90c}.
        \item Jan 2027: Benchmark port vs spoofing, run profiling \& optimise fixes. Check \& resolve regressions as per ROCm test suite.
        \item Feb 2027: Validate approach on \texttt{gfx902}
        \item Mar 2027: Set up CI/CID for all unsupported AMD GCN GPUs.
        \item Apr-Jul 2027: Contribute upstream to Mojo, LLVM, ROCm, PyTorch, and vLLM with RDNA 4 development platform.
    \end{itemize}
\end{frame}

\begin{frame}{Literature review -- ROCm software stack}
\begin{itemize}
    \item ROCm Documentation (\href{https://rocm.docs.amd.com/}{\texttt{rocm.docs.amd.com/...}})
    \item HIP Programming Model (\href{https://rocm.docs.amd.com/projects/HIP/en/latest/understand/programming_model.html}{\texttt{rocm.docs.amd.com/.../programming\_model}})
    \item HIP Programming Patterns (\href{https://rocm.docs.amd.com/projects/HIP/en/latest/tutorial/programming-patterns.html}{\texttt{rocm.docs.amd.com/.../programming-patterns}})
    \item HIP SAXPY Tutorial (\href{https://rocm.docs.amd.com/projects/HIP/en/latest/tutorial/saxpy.html}{\texttt{rocm.docs.amd.com/.../saxpy}})
    \item HIPIFY (\href{https://github.com/ROCm/HIPIFY}{\texttt{github.com/ROCm/HIPIFY}})
    \item ROCm Examples (\href{https://github.com/ROCm/rocm-examples}{\texttt{github.com/ROCm/rocm-examples}})
    \item ROCm Libraries (\href{https://github.com/ROCm/rocm-libraries}{\texttt{github.com/ROCm/rocm-libraries}})
    \item rocBLAS (\href{https://github.com/ROCm/rocBLAS}{\texttt{github.com/ROCm/rocBLAS}})
    \item hipBLAS (\href{https://github.com/ROCm/hipBLAS}{\texttt{github.com/ROCm/hipBLAS}})
    \item ROCm Platform (\href{https://github.com/ROCm/ROCm}{\texttt{github.com/ROCm/ROCm}})
    \item ROCm LLVM Fork (\href{https://github.com/ROCm/llvm-project}{\texttt{github.com/ROCm/llvm-project}})
\end{itemize}
\end{frame}


\begin{frame}{Literature review -- Compiler, build system and upstream development}
\begin{itemize}
    \item ROCm Libraries PR \#8139 (\href{https://github.com/ROCm/rocm-libraries/pull/8139}{\texttt{github.com/ROCm/.../pull/8139}})
    \item TheRock Issue \#3818 (\href{https://github.com/ROCm/TheRock/issues/3818}{\texttt{github.com/ROCm/.../issues/3818}})
    \item TheRock Build System (\href{https://github.com/ROCm/TheRock}{\texttt{github.com/ROCm/TheRock}})
    \item ROCm GitHub Organisation (\href{https://github.com/ROCm}{\texttt{github.com/ROCm}})
    \item RadeonOpenCompute (\href{https://github.com/RadeonOpenCompute}{\texttt{github.com/RadeonOpenCompute}})
    \item LLVM AMDGPU Backend (\href{https://github.com/llvm/llvm-project/tree/main/llvm/lib/Target/AMDGPU}{\texttt{github.com/llvm/.../AMDGPU}})
    \item LLVM Clang (\href{https://github.com/llvm/llvm-project/tree/main/clang}{\texttt{github.com/llvm/.../clang}})
    \item ROCm LLVM Toolchain (\href{https://github.com/ROCm/llvm-project}{\texttt{github.com/ROCm/llvm-project}})
    \item AMD GCN ISA Reference (\href{https://gpuopen.com/learn/amd-gcn-assembly-cross-lane-operations/}{\texttt{gpuopen.com/.../gcn-assembly}})
    \item GPUOpen (\href{https://gpuopen.com/}{\texttt{gpuopen.com}})
\end{itemize}
\end{frame}


\begin{frame}{Literature review -- GPU architecture and supporting references}
\begin{itemize}
    \item Graphics Processing Unit (\href{https://en.wikipedia.org/wiki/Graphics_processing_unit}{\texttt{wikipedia.org/.../Graphics\_processing\_unit}})
    \item Geometry Engine (\href{https://en.wikipedia.org/wiki/Geometry_Engine}{\texttt{wikipedia.org/.../Geometry\_Engine}})
    \item Transformation Matrix (\href{https://en.wikipedia.org/wiki/Transformation_matrix}{\texttt{wikipedia.org/.../Transformation\_matrix}})
    \item GPUOpen Technical Articles (\href{https://gpuopen.com/learn/}{\texttt{gpuopen.com/learn}})
    \item Project Research Notes (\href{https://docs.google.com/document/d/1dZ3vF8GE8_gx6t152sOaUVEPq0ybmailxvu3uk89_is/edit}{\texttt{docs.google.com/...}})
    \item Azu (\href{https://github.com/Netherquark/azu}{\texttt{github.com/Netherquark/azu}})
    \item PuneFOSS CFP (\href{https://fossunited.org/c/pune/punefoss4/cfp/2lhj0bntp3}{\texttt{fossunited.org/.../2lhj0bntp3}})
    \item ROCm Examples Repository (\href{https://github.com/ROCm/rocm-examples}{\texttt{github.com/ROCm/rocm-examples}})
    \item RadeonOpenCompute Organisation (\href{https://github.com/RadeonOpenCompute}{\texttt{github.com/RadeonOpenCompute}})
    \item HIPIFY Repository (\href{https://github.com/ROCm/HIPIFY}{\texttt{github.com/ROCm/HIPIFY}})
\end{itemize}
\end{frame}

\begin{frame}{Gap in the research, technology, and methodology}
\vspace{-0.5cm}
\begin{itemize}
        \item No documented end-to-end workflow for recovering unsupported ROCm targets.
        \item Packaged ROCm components omitted \texttt{gfx90c}-specific code generation and runtime data.
        \item Recovery knowledge was fragmented across documentation, issue trackers, pull requests, and community discussions.
        \item No methodology existed to validate recovery, benchmark correctness, and reproduce results across GPU generations.
    \end{itemize}
\end{frame}

\begin{frame}{Gap in the research, technology, and methodology}
\vspace{-1.5cm}
\hspace{-3cm}
\begin{figure}
\includegraphics[width=1.1\linewidth]{sankey.png}
\end{figure}
\end{frame}


\begin{frame}{Description of the proposed solution}
\begin{itemize}
        \item Establish a staged recovery workflow from native binaries to complete source builds.
        \item Validate every layer using HIP applications, \texttt{rocm-examples}, and PyTorch workloads.
        \item Isolate compiler, runtime, library, packaging, and hardware-specific failures independently.
        \item Produce a reproducible workflow transferable to other unsupported AMD GCN GPUs.
\end{itemize}
\end{frame}

\begin{frame}{Requirement analysis}
\begin{figure}
\includegraphics[width=\linewidth]{system-info.png}
\end{figure}
\end{frame}

\begin{frame}{Technology Stack}

\centering
\setlength{\tabcolsep}{7pt}

%---------------- Row 1 ----------------%
\begin{tabular}{ccccc}
{\parbox[c][1.55cm][c]{1.55cm}{\centering 
\includegraphics[width=\linewidth]{fedora-logo.png}}} &
{\parbox[c][1.55cm][c]{1.55cm}{\centering 
\includegraphics[width=\linewidth]{rocm-logo.png}}} &
{\parbox[c][1.55cm][c]{1.55cm}{\centering 
\includegraphics[width=\linewidth]{llvm-logo.png}}} &
{\parbox[c][1.55cm][c]{1.55cm}{\centering 
\includegraphics[width=0.9\linewidth]{cpp-logo.png}}} &
{\parbox[c][1.55cm][c]{1.55cm}{\centering 
\includegraphics[width=\linewidth]{git-logo.png}}}
\end{tabular}

\vspace{0.5cm}

%---------------- Row 2 ----------------%
\begin{tabular}{ccccc}
{\parbox[c][1.55cm][c]{1.55cm}{\centering 
\includegraphics[width=\linewidth]{bash-logo.png}}} &
{\parbox[c][1.55cm][c]{1.55cm}{\centering 
\includegraphics[width=0.7\linewidth]{pytorch-logo.png}}} &
{\parbox[c][1.55cm][c]{1.55cm}{\centering 
\includegraphics[width=\linewidth]{mojo-logo.png}}} &
{\parbox[c][1.55cm][c]{1.55cm}{\centering 
\includegraphics[width=\linewidth]{python-logo.png}}} &
{\parbox[c][1.55cm][c]{1.55cm}{\centering 
\includegraphics[width=\linewidth]{geekbench-logo.png}}}
\end{tabular}
\end{frame}

\begin{frame}{Native ROCm baseline}
\vspace{-0.5cm}
\begin{itemize}
    \item Installed ROCm 6.4.4 on Fedora 43 for \texttt{gfx90c}.
    \item Verified HIP compilation and execution using native ROCm binaries.
    \item Identified PyTorch failures caused by missing \texttt{gfx90c} rocBLAS Tensile data.
    \item Evaluated HSA architecture spoofing as an initial compatibility baseline.
    \item Ran Fedora provided \texttt{rocm-examples} as a baseline.
\end{itemize}
\end{frame}


\begin{frame}{Source recovery}
\vspace{-0.5cm}
\begin{itemize}
    \item Built \texttt{rocm-examples} from source using the ROCm 6.4.4 toolchain.
    \item Investigated missing \texttt{main\_gfx90c.s} assembly generation.
    \item Extended the \texttt{assembly\_to\_executable} workflow for \texttt{gfx90c}.
    \item Traced compilation from HIP source through LLVM IR to AMDGPU assembly.
\end{itemize}
\end{frame}


\begin{frame}{Failure analysis}
\vspace{-0.5cm}
\begin{itemize}
    \item Classified failures into packaging, metadata, missing Tensile data, and hardware limitations.
    \item Compared behaviour across native, spoofed, and source-built environments.
    \item Validated each recovery stage using HIP applications and \texttt{rocm-examples}.
    \item Recorded reproducible failure cases for subsequent recovery work.
\end{itemize}
\end{frame}


\begin{frame}{Benchmark evaluation}
\vspace{-0.5cm}
\begin{itemize}
    \item Evaluated memory bandwidth using CLBlast and clpeak.
    \item Measured matrix multiplication performance using rocBLAS benchmark clients.
    \item Compared application behaviour under native and spoofed execution.
    \item Used benchmark results to validate recovered functionality and identify remaining regressions.
\end{itemize}
\end{frame}



\begin{frame}{Benchmark evaluation}
\vspace{-3cm}
\hspace{-5cm}
\begin{figure}[htbp]
    \includegraphics[width=\linewidth]{spoofing_summary.png}
\end{figure}
\end{frame}

\begin{frame}{Benchmark evaluation}
\begin{figure}[htbp]
    \includegraphics[width=\linewidth]{spoofing-performance.png}
\end{figure}
\end{frame}

\begin{frame}{Benchmark evaluation}
\begin{figure}[htbp]
    \includegraphics[width=\linewidth]{hip-basic-progress.png}
\end{figure}
\end{frame}

\begin{frame}{Benchmark evaluation}
\begin{figure}[htbp]
    \includegraphics[width=\linewidth]{overall_progress.png}
\end{figure}
\end{frame}

\begin{frame}{Intended project outcome}
\vspace{-0.5cm}
\begin{itemize}
    \item A reproducible workflow for recovering unsupported AMD GCN GPUs under ROCm.
    \item Complete \texttt{gfx90c} enablement across the ROCm software stack.
    \item Validation of the recovery methodology on \texttt{gfx902}.
    \item Benchmark-backed evaluation of native, spoofed, and recovered software stacks.
    \item Documentation suitable for upstream contribution and long-term maintenance.
\end{itemize}
\end{frame}


\begin{frame}{Conclusion}
\vspace{-0.5cm}
\begin{itemize}
    \item Native ROCm functionality on unsupported hardware was successfully characterised.
    \item Compiler, runtime, packaging, and library limitations were systematically isolated.
    \item A staged recovery methodology was established for unsupported ROCm targets.
    \item The proposed workflow provides a foundation for future upstream enablement of legacy AMD GPUs.
\end{itemize}
\end{frame}


\begin{frame}{References}
\scriptsize
\begin{itemize}
    \item AMD ROCm Documentation\\
    \url{https://rocm.docs.amd.com/}

    \item AMD ROCm Programming Guide\\
    \url{https://rocm.docs.amd.com/projects/HIP/en/latest/}

    \item HIP Programming Model\\
    \url{https://rocm.docs.amd.com/projects/HIP/en/latest/understand/programming_model.html}

    \item HIP Programming Patterns\\
    \url{https://rocm.docs.amd.com/projects/HIP/en/latest/tutorial/programming-patterns.html}

    \item HIP SAXPY Tutorial\\
    \url{https://rocm.docs.amd.com/projects/HIP/en/latest/tutorial/saxpy.html}

    \item ROCm GitHub Organisation\\
    \url{https://github.com/ROCm}

    \item ROCm Platform Repository\\
    \url{https://github.com/ROCm/ROCm}

    \item ROCm Libraries Repository\\
    \url{https://github.com/ROCm/rocm-libraries}

    \item ROCm Libraries Pull Request \#8139\\
    \url{https://github.com/ROCm/rocm-libraries/pull/8139}

    \item ROCm Examples Repository\\
    \url{https://github.com/ROCm/rocm-examples}

    
\end{itemize}
\end{frame}

\begin{frame}{References}
    \scriptsize
    \begin{itemize}
        \item rocBLAS Repository\\
    \url{https://github.com/ROCm/rocBLAS}

    \item hipBLAS Repository\\
    \url{https://github.com/ROCm/hipBLAS}

    \item TheRock Repository\\
    \url{https://github.com/ROCm/TheRock}

    \item TheRock Issue \#3818\\
    \url{https://github.com/ROCm/TheRock/issues/3818}

    \item ROCm LLVM Fork\\
    \url{https://github.com/ROCm/llvm-project}

    \item LLVM Project\\
    \url{https://github.com/llvm/llvm-project}

    \item LLVM AMDGPU Backend\\
    \url{https://github.com/llvm/llvm-project/tree/main/llvm/lib/Target/AMDGPU}

    \item GPUOpen\\
    \url{https://gpuopen.com/}

    \item AMD GCN Assembly Article\\
    \url{https://gpuopen.com/learn/amd-gcn-assembly-cross-lane-operations/}

    \item RadeonOpenCompute\\
    \url{https://github.com/RadeonOpenCompute}

    
    \end{itemize}
\end{frame}

\begin{frame}{References}
   \scriptsize
    \begin{itemize}
    \item HIPIFY Repository\\
    \url{https://github.com/ROCm/HIPIFY}
    \item Graphics Processing Unit\\
    \url{https://en.wikipedia.org/wiki/Graphics_processing_unit}

    \item Geometry Engine\\
    \url{https://en.wikipedia.org/wiki/Geometry_Engine}

    \item Transformation Matrix\\
    \url{https://en.wikipedia.org/wiki/Transformation_matrix}
        \item Azu\\
    \url{https://github.com/Netherquark/azu}

    \item PuneFOSS 2025 CFP\\
    \url{https://fossunited.org/c/pune/punefoss4/cfp/2lhj0bntp3}

    \item AMD vis-\`a-vis GPUs Blog Series (\#1--\#8)

    \item Progress Report No.~1

    \item Progress Report No.~2
    \end{itemize}
\end{frame}

\begin{frame}{31 July 2026 CA-1 Presentation}
    \begin{figure}[htbp]
        \includegraphics[width=\linewidth]{meeting.png}
    \end{figure}
\end{frame}

%\begin{frame}{Project outcome}
%\begin{columns}[T,onlytextwidth]
%    \column{0.55\textwidth}
%    \begin{itemize}
%        \item gfx90c is now a studied, testable target rather than an unknown.
%        \item The runtime path works for basic HIP code.
%        \item The remaining blockers sit in packaging, kernels, and generated artifacts.
%        \item gfx902 is ready for the same validation sequence.
%    \end{itemize}
%    \column{0.45\textwidth}
%    \figph{4.0cm}{\linewidth}{Outcome dashboard placeholder}
%\end{columns}
%\end{frame}
%
%
%\begin{frame}{Conclusion}
%\begin{columns}[T,onlytextwidth]
%    \column{0.55\textwidth}
%    \begin{itemize}
%        \item The project moved from a broad idea to a reproducible recovery method.
%        \item The strongest evidence points to software and packaging gaps.
%        \item The next stage is to finish the gfx90c port and then harden the pipeline.
%        \item The later roadmap expands into direct upstream work after new hardware arrives.
%    \end{itemize}
%    \column{0.45\textwidth}
%    \figph{4.0cm}{\linewidth}{Roadmap end-state placeholder}
%\end{columns}
%\end{frame}
%
%
%\begin{frame}{References}
%\small
%\begin{itemize}
%    \item AMD vis-\`a-vis GPUs \#1 - \#8.
%    \item Progress Report No \#1 and Progress Report No \#2.
%    \item AMD ROCm documentation, LLVM source history, and the Vega ISA reference guide.
%    \item Fedora packaging notes and issue threads related to gfx90c support.
%\end{itemize}
%\end{frame}

\end{document}
